% Part: turing-machines
% Chapter: machines-computations
% Section: combining-machines

\documentclass[../../../include/open-logic-section]{subfiles}

\begin{document}

\olfileid{tur}{mac}{cmb}
\olsection{ट्यूरिंग यंत्रांचे संयोजन}

\begin{explain}
आतापर्यंत पाहिलेली ट्यूरिंग यंत्रांची उदाहरणे बरीच सोपी आहेत.
पण आधुनिक कार्यक्रमणभाषेत सोडवता येणारी कोणतीही समस्या ट्यूरिंग
यंत्रांनीही सोडवता येते. अधिक गुंतागुंतीची ट्यूरिंग यंत्रे बनवताना
ती एकत्र जोडता येतात हे समजून घेणे महत्त्वाचे आहे. मग प्रक्रिया
सोप्या भागांत विभागून अधिक गुंतागुंतीच्या समस्या सोडवणारी यंत्रे
बनवता येतात. गुंतागुंतीची समस्या घटक भागांत विभागण्याचा नैसर्गिक
मार्ग सापडला, तर अनेक सोपी ट्यूरिंग यंत्रे तयार करून त्यांचे
एकाच यंत्रात संयोजन करता येते. असे यंत्र अनेक टप्प्यांत समस्या
सोडवते. पुढील विभागातील थांबण्याची समस्या हाताळताना हे विशेष
महत्त्वाचे ठरेल.

$M = \tuple{Q, \Sigma, q_0, \delta}$ आणि
~$M' = \tuple{Q', \Sigma', q_0', \delta'}$ ही ट्यूरिंग यंत्रे
कशी जोडायची? $M$~थांबल्यानंतर फितीचे जे स्थितिवर्णन असते तेच
यंत्र~$M'$ च्या धावेचे प्रारंभिक स्थितिवर्णन म्हणून वापरू.
असे करणारे एकच ट्यूरिंग यंत्र $M \frown M'$ मिळवण्यासाठी पुढील
पायऱ्या घ्या:
\begin{enumerate}
    \item $M'$ च्या~$Q'$ मधील सर्व अवस्थांचे क्रमांक (किंवा नावे)
    बदला, म्हणजे $M$ आणि~$M'$ यांची कोणतीही अवस्था समान राहणार
    नाही ($Q \cap Q' = \emptyset$).
    \item $M \frown M'$ च्या अवस्था $Q \cup Q'$ असतील.
    \item फितीची वर्णमाला $\Sigma \cup \Sigma'$ असेल.
    \item प्रारंभिक अवस्था~$q_0$ असेल.
    \item संक्रमण फलन पुढील $\delta''$ असेल:
    \[\delta''(q,\sigma) =
    \begin{cases}
      \delta(q,\sigma) & \text{$q \in Q$ आणि $\delta(q,\sigma)$ परिभाषित असेल तर}\\
      \delta'(q,\sigma) & \text{$q \in Q'$ असेल तर}\\
      \tuple{q_0', \sigma, \TMstay} & \text{$q \in Q$ असेल आणि
      $\delta(q,\sigma)$ अपरिभाषित असेल तर}
    \end{cases}\]
\end{enumerate}
%READERNOTE{OLTUR-007 — गोठवलेल्या इंग्रजीतील संक्रमण फलनाच्या पहिल्या शाखेत फक्त अवस्था Q मध्ये असणे ही अट आहे; त्यामुळे तिसऱ्या, अपरिभाषित-संक्रमण शाखेशी ती छेदते. येथे पहिली शाखा δ परिभाषित असतानाच लागू होईल असे स्पष्ट केले आहे. सर्व चिन्हे आणि इतर शाखा जतन आहेत.}
मिळालेले यंत्र $q \in Q$ अवस्थेत असताना~$M$ च्या सूचना,
तर $q \in Q'$ अवस्थेत असताना~$M'$ च्या सूचना वापरते.
$q \in Q$ अवस्थेत~$\sigma$ चिन्ह वाचताना $M$ मध्ये त्यासाठी
संक्रमण नसेल (म्हणजे $M$ थांबले असते), तर ते~$M'$ च्या
प्रारंभिक अवस्थेत जाते; फितीवरील मजकूर आणि शीर्षाची जागा
तशीच ठेवते.

यंत्र~$M$ शिस्तबद्ध नसेल तर $M$~थांबताना फितीचे शीर्ष कुठे
आहे हे आपल्याला माहीत नसते. म्हणून~$M$ च्या अंतिम
स्थितिवर्णनात शीर्ष घर~$1$ वाचत असेलच असे नाही. यंत्रे
जोडताना हे लक्षात ठेवणे महत्त्वाचे आहे.
\end{explain}

\begin{ex}
\emph{यंत्रांचे संयोजन:} सुरुवातीला $n$ आणि~$m$ लांबीचे
$\TMstroke$ चिन्हांचे दोन खंड आदान म्हणून घेऊन शेवटी
फितीवर $2(m+n)$ $\TMstroke$ चिन्हांचा एकच खंड ठेवून थांबणारे
यंत्र रचू. यासाठी परिचित अशी बेरीज-यंत्र आणि दुप्पट करणारे
यंत्र ही दोन यंत्रे जोडता येतील. आधी बेरीज-यंत्राचा अवस्था-आलेख
काढू.
\[
\begin{tikzpicture}[->,>=stealth',shorten >=1pt,auto,node distance=2.8cm,
                    semithick]
  \tikzstyle{every state}=[fill=none,draw=black,text=black]

  \node[initial,state] (A)              {$q_0$};
  \node[state]         (B) [right of=A] {$q_1$};
  \node[state]         (C) [right of=B] {$q_2$};

  \path (A) edge node {\TMtrans{\TMblank}{\TMstroke}{\TMstay}} (B)
            edge [loop above] node {\TMtrans{\TMstroke}{\TMstroke}{\TMright}} (B)
        (B) edge [loop above] node {\TMtrans{\TMstroke}{\TMstroke}{\TMright}} (B)
            edge node {\TMtrans{\TMblank}{\TMblank}{\TMleft}} (C)
        (C) edge [loop above] node {\TMtrans{\TMstroke}{\TMblank}{\TMstay}} (C);
\end{tikzpicture}
\]
प्रदान दुप्पट करण्यासाठी अवस्था~$q_2$ मध्ये थांबण्याऐवजी
यंत्र पुढे चालवायचे आहे. दुप्पट करणारे यंत्र आदानातील पहिली
रेघ पुसते आणि स्वतंत्र प्रदानात दोन रेघा लिहिते, हे आठवा.
बेरीज-यंत्राच्या प्रदानातील पहिली रेघ शीर्ष वाचत असेल याची
खात्री करणारी सूचना जोडू.
\[
\begin{tikzpicture}[->,>=stealth',shorten >=1pt,auto,node distance=2.8cm,
                    semithick]
  \tikzstyle{every state}=[fill=none,draw=black,text=black]

  \node[initial,state] (A)              {$q_0$};
  \node[state]         (B) [right of=A] {$q_1$};
  \node[state]         (C) [right of=B] {$q_2$};
  \node[state]         (D) [below left of=C] {$q_3$};
  \node[state]         (E) [below left of=D] {$q_4$};

  \path (A) edge node {\TMtrans{\TMblank}{\TMstroke}{\TMstay}} (B)
            edge [loop above] node {\TMtrans{\TMstroke}{\TMstroke}{\TMright}} (B)
        (B) edge [loop above] node {\TMtrans{\TMstroke}{\TMstroke}{\TMright}} (B)
            edge node {\TMtrans{\TMblank}{\TMblank}{\TMleft}} (C)
        (C) edge node {\TMtrans{\TMstroke}{\TMblank}{\TMleft}} (D)
        (D) edge [loop left] node {\TMtrans{\TMstroke}{\TMstroke}{\TMleft}} (D)
            edge node[left, xshift=-2mm] {\TMtrans{\TMendtape}{\TMendtape}{\TMright}} (E);
\end{tikzpicture}
\]
आता आदान दुप्पट करणे सोपे आहे: दुप्पट करणारे यंत्र
अवस्था~$q_4$ ला जोडायचे. त्यासाठी त्या यंत्राच्या अवस्थांना
पुन्हा नावे द्यावी लागतील, म्हणजे त्या~$q_0$ ऐवजी~$q_4$
पासून सुरू होतील. अशा रीतीने दोन प्रारंभिक अवस्था राहणार नाहीत.
अंतिम अवस्था-आलेख \olref{fig:combined} मध्ये दाखवल्याप्रमाणे
असेल.
\begin{figure}
\[
\begin{tikzpicture}[->,>=stealth',shorten >=1pt,auto,node distance=2.8cm,
                    semithick]
  \tikzstyle{every state}=[fill=none,draw=black,text=black]
  \node[initial,state] (A)              {$q_0$};
  \node[state]         (B) [right of=A] {$q_1$};
  \node[state]         (C) [right of=B] {$q_2$};
  \node[state]         (D) [below left of=C] {$q_3$};
  \node[state]         (E) [below left of=D] {$q_4$};
  \node[state]         (2) [right of=E] {$q_5$};
  \node[state]         (3) [right of=2] {$q_6$};
  \node[state]         (4) [below of=3] {$q_7$};
  \node[state]         (5) [left of=4]  {$q_8$};
  \node[state]         (6) [left of=5]  {$q_9$};

  \path (A) edge node {\TMtrans{\TMblank}{\TMstroke}{\TMstay}} (B)
            edge [loop above] node {\TMtrans{\TMstroke}{\TMstroke}{\TMright}} (B)
        (B) edge [loop above] node {\TMtrans{\TMstroke}{\TMstroke}{\TMright}} (B)
            edge node {\TMtrans{\TMblank}{\TMblank}{\TMleft}} (C)
        (C) edge node {\TMtrans{\TMstroke}{\TMblank}{\TMleft}} (D)
        (D) edge [loop left] node {\TMtrans{\TMstroke}{\TMstroke}{\TMleft}} (D)
            edge node[left, xshift=-2mm] {\TMtrans{\TMendtape}{\TMendtape}{\TMright}} (E)
    (E) edge node {\TMtrans{\TMstroke}{\TMblank}{\TMright}} (2)
    (2) edge [loop above] node {\TMtrans{\TMstroke}{\TMstroke}{\TMright}} (2)
      edge node {\TMtrans{\TMblank}{\TMblank}{\TMright}} (3)
    (3) edge [loop above] node {\TMtrans{\TMstroke}{\TMstroke}{\TMright}} (3)
        edge node {\TMtrans{\TMblank}{\TMstroke}{\TMright}} (4)
    (4) edge [loop below] node {\TMtrans{\TMblank}{\TMstroke}{\TMleft}} (4)
        edge node {\TMtrans{\TMstroke}{\TMstroke}{\TMleft}} (5)
    (5) edge [loop below]  node {\TMtrans{\TMstroke}{\TMstroke}{\TMleft}} (5)
        edge              node {\TMtrans{\TMblank}{\TMblank}{\TMleft}} (6)
    (6) edge [loop below] node {\TMtrans{\TMstroke}{\TMstroke}{\TMleft}} (6)
        edge              node {\TMtrans{\TMblank}{\TMblank}{\TMright}} (E);
\end{tikzpicture}
\]
\caption{बेरीज-यंत्र आणि दुप्पट करणारे यंत्र यांचे संयोजन}
\ollabel{fig:combined}
\end{figure}
\end{ex}

\begin{prop}
    $M$ आणि $M'$ ही शिस्तबद्ध यंत्रे अनुक्रमे $f\colon
    \Nat^k \to \Nat$ आणि $f'\colon \Nat \to \Nat$ ही फलने
    संगणित करत असतील, तर $M \frown M'$ हेही शिस्तबद्ध असून
    ~$\comp{f}{f'}$ संगणित करते.
\end{prop}

\begin{proof}
    $M$ शिस्तबद्ध असल्यामुळे ते
    प्रदान~$f(n_1,\dots,n_k) = m$ देऊन थांबते तेव्हा शीर्ष
    घर~$1$ वाचत असते. आता यंत्र~$M'$ च्या प्रारंभिक अवस्थेत
    गेल्यावर प्रदान $f'(m)$ देऊन यंत्र पुन्हा घर~$1$ वरच थांबते.
    \olref[dis]{defn:disciplined} मधील इतर अटीही पूर्ण होतात.
\end{proof}

\begin{prob}
    \cref{tur:mac:dis:prob:disc-succ} मधील यंत्र~$M$ घेऊन
    $M \frown M$ रचा आणि $f(x) = x+2$ संगणित करणारे
    शिस्तबद्ध ट्यूरिंग यंत्र द्या.
\end{prob}
\end{document}
